\documentclass[a5paper,10pt]{article}

% https://github.com/InDevRus/filippov-solutions
% Нумерация страниц
\pagenumbering{gobble}

% Настройка полей
\usepackage{geometry}
\geometry{left=1cm}
\geometry{right=1cm}
\geometry{top=1cm}
\geometry{bottom=1.5cm}

% Вставка таблицы
\usepackage{array}
\usepackage{tabularx}

% Инструменты выравнивания формул
\usepackage{amsmath}
\usepackage{mathtools}

% Диакритика в математике
\usepackage{amsfonts}

% Вычисления размеров на ходу
\usepackage{calc}

% Удобные записи производных
\usepackage[italic=true]{derivative}

% Настройки сносок
\usepackage{enotez}

% Длинные знаки векторов
\usepackage{anyfontsize}
\usepackage[b]{esvect}

% Вставка картинок
\usepackage{graphicx}

% Вставка красной строки
\usepackage{indentfirst}
\setlength{\parindent}{1em}

% Условия в командах
\usepackage{xifthen}

% Луа-скрипты
\usepackage{luacode}

% Настройка команд отдельно в математическом и текстовом режимах
\usepackage{mathcommand}

% Для повторения LaTeX кода
\usepackage{multido}

% Конфигурация отступов формул
\delimitershortfall-1sp
\usepackage{mleftright}
\mleftright

% Растягивание объектов
\usepackage{relsize}

% Обтекание текстом
\usepackage{wrapfig}
\usepackage{subcaption}
\usepackage{floatflt}

% Поддержка ввода кириллицы
\usepackage[english, russian]{babel}

% Настройка корневой позиции для компилятора
\usepackage{import}
\providecommand*{\ProjectRootPath}{..}

\import{\ProjectRootPath/preambles/macros/}{theorems}

\import{\ProjectRootPath/preambles/fonts}{font_setup}

% Знак градуса
\usepackage{siunitx}

\import{\ProjectRootPath/preambles/macros/}{operators}
\import{\ProjectRootPath/preambles/macros/}{redefinitions}

\import{\ProjectRootPath/preambles/fonts}{letters_setup}
\import{\ProjectRootPath/preambles/fonts/adjustments}{custom_superscripts}

\import{\ProjectRootPath/preambles/custom}{formatting}
\import{\ProjectRootPath/preambles/custom}{frames}
\import{\ProjectRootPath/preambles/custom}{list_setup}

% TODO: перевести команды на современный стиль объявления

\begin{document}

\begin{framed}
    Для угла $\phi = \ang{90}$ имеем $\tg{\beta} = \tg\br{\alpha \pm \phi} =\tg\br{\alpha \pm \ang{90}} = -\ctg \alpha$.
\end{framed}

$\pow{x}{2} + \pow{C}{2} = 2Cy$.
$2x = 2Cy\'$.
$C = \dfrac {x} {y\'}$.
$\pow{x}{2} + \br{\dfrac {x}{y\'}}^2 = 2\br{\dfrac {x} {y\'}} y$.
$\br{y\'}^2 x - 2yy\' + x = 0$.
$x \pow{\tg}{2} \alpha - 2y\tg\alpha + x = 0$.
Подставив $\tg\alpha = -\dfrac {1} {\tg\beta}$, получаем
$\dfrac {x} {\pow{\tg}{2} \beta} + \dfrac {2y} {\tg\beta} + x = 0$, так что
$x \pow{\tg}{2} \beta + 2y\tg\beta + x = 0$.
Окончательно, уравнение изогональных (ортогональных) траекторий имеет вид 
$x\br{z\'}^2 + 2zz\' + x = 0$.

\begin{remark}
    Уравнение $x\br{z\'}^2 + 2zz\' + x = 0$ имеет общее решение в параметрической форме
    $\tsys{
        & x\br{t} = 2C {t} {\br{3\pow{t}{2} + 1}}^{\mathlarger{-\frac {2} {3}}} \\& z\br{t} = -C \br{\pow{t}{2} + 1} {\br{3\pow{t}{2} + 1}}^{\mathlarger{-\frac {2} {3}}}
    }$.
\end{remark}

\end{document}

